\section{6. Discussion}\label{discussion}

\subsection{6.1 What is supported}\label{what-is-supported}

The supported result is narrow: for these two reconstructed published
curves, fitted intercepts depend materially on the declared functional
form and on lower fit-domain membership. This dependence remains after
boundary-limited fits are excluded and cannot be attributed to the
isolated two-point plateau.

\subsection{6.2 Why the source model
matters}\label{why-the-source-model-matters}

The plotted \(\alpha(E)\) curves inherit assumptions from the original
ellipsometric inversion. Thickness, interface roughness, dielectric
parameterization, spectral resolution, and covariance can affect the
recovered optical constants before the present extraction model is
applied \citep{Ibrahim1971, BuAbbud1981, GarciaCaurel2013}. The present
analysis does not propagate those upstream effects. It should therefore
be interpreted as a lower-level audit of downstream reporting
conventions, not as a total uncertainty budget for the specimens.

\subsection{6.3 Why the model family is not a probability
distribution}\label{why-the-model-family-is-not-a-probability-distribution}

The extraction models were chosen because they represent historically
used or source-supported conventions. They were not sampled from a prior
distribution, and their range is not a 68\%, 95\%, or worst-case
physical confidence interval. The span answers a deterministic question:
how far do the reported intercepts move among the declared models?

\subsection{6.4 Model adequacy and physical
completeness}\label{model-adequacy-and-physical-completeness}

A small residual does not prove that a fitted intercept equals the
latent gap. Conversely, a broad model family should not be allowed to
inflate a headline range with plainly incompatible forms. Residual
diagnostics are therefore reported explicitly, while formal stochastic
model selection is avoided. The residual-envelope screen is
intentionally demoted because its retained span changes under stricter
coverage thresholds; it is evidence that some models are visibly
incompatible, not a unique model-selection rule. More complete HgCdTe
absorption models incorporating nonparabolic conduction and light-hole
bands, heavy-hole contributions, and Urbach tails remain physically
important \citep{Finkman1984, Chang2006, Chang2007}. They require
additional source-complete parameters and are not approximated by
undocumented nuisance choices here.

\subsection{6.5 Future validation}\label{future-validation}

The strongest next step is not another equation ranking from plotted
curves. It is analysis of native spectral data from multiple specimens
and laboratories with:

\begin{itemize}
\tightlist
\item
  independent composition metrology and uncertainty;
\item
  carrier type and density;
\item
  thickness and interface-model covariance;
\item
  native ellipsometric observables or instrument-export absorption data;
\item
  predeclared extraction models and fit domains;
\item
  same-specimen comparison with an independent gap-sensitive modality
  where feasible.
\end{itemize}

A paired factorial design can be useful, but it is outside the
evidentiary core of this paper and is not presented as a main result.
